
\section{\trFloer{}算子的（非线性）\trEllipticRegularity：证明}
\label{sec:nonlinear-proofs}

\subsection{如何推出命题~12.1.4……}
\label{subsec:deduce-prop-1214}

……由定理~12.1.5。我们假设这条定理已经获证。若对某个 $k\ge 1$ 有 $u\in W^{k,p}_{\mathrm{loc}}(\RR\times S^{1};\RR^{2n})$，那么这条定理便蕴含 $u\in W^{k+1,p}_{\mathrm{loc}}(\RR\times S^{1};\RR^{2n})$，因为 $\widetilde{J}$ 与 $h$ 同 $u$ 具有相同的\trRegularity。于是，对每个 $k$ 都有 $u\in W^{k,p}_{\mathrm{loc}}(\RR\times S^{1};\RR^{2n})$，并且由
\[
\bigcap W^{k,p}_{\mathrm{loc}}\subset\cC^{\infty}\text{，}
\]
这一事实——它由\trRellich{}定理（\cref{thm:sobolev-estimates}）可得——可知 $u$ 是 $\cC^{\infty}$ 类的。

为了证明\cref{thm:loc-topologies-coincide}的第二条断言，我们先来证明下述结果。

\begin{corollary}\label{thm:smooth-solution-wkp-estimate}
设 $u\in\cC^{\infty}(\RR\times S^{1};\RR^{2n})$ 是\trFloerEquation
\[
\frac{\partial u}{\partial s}+\widetilde{J}\frac{\partial u}{\partial t}+h=0
\]
的一个解（其中 $\widetilde{J}\in\cC^{\infty}\bigl(\RR\times S^{1};\operatorname{End}(\RR^{2n})\bigr)$，$\widetilde{J}^{2}=-\mathrm{Id}$，$h\in\cC^{\infty}(\RR\times S^{1};\RR^{2n})$）。设 $V\subset U\subset\RR\times S^{1}$ 是两个相对紧的\trOpen{}子集，并设 $p>2$、$k\ge 1$ 与 $c_{0}>0$ 满足
\[
\|\widetilde{J}\|_{W^{k,p}(U)}\le c_{0}\text{。}
\]
那么存在一个常数 $c$（它依赖于 $V$、$U$、$c_{0}$、$k$ 与 $p$），使得
\[
\|u\|_{W^{k+1,p}(V)}\le c\left(\|h\|_{W^{k,p}(U)}+\|u\|_{W^{k,p}(U)}\right)\text{。}
\]
\end{corollary}

这正是\cref{thm:generalized-regularity-estimate}的那条不等式，对每个 $k\in\NN$ 皆是。

\begin{proof}[\cref{thm:smooth-solution-wkp-estimate}的证明]
由于\cref{thm:generalized-regularity-estimate}的条件对每个 $k$ 都满足，我们便有所要的那条不等式。
\end{proof}

我们首先假设 $(u_{n})$ 是一个在 $W^{1,p}_{\mathrm{loc}}(\RR\times S^{1};\RR^{2n})$ 中\trConvergence{}到 $u$ 的序列，其中每一个 $u_{n}$ 连同 $u$ 自己，都由上文知是\trFloerEquation{}的 $\cC^{\infty}$ 类解。令 $\widetilde{J}_{n}(s,t)=J(u_{n}(s,t))$，并令 $h_{n}(s,t)=\operatorname{grad}H_{t}(u_{n}(s,t))$，于是（按 $\widetilde{J}$ 与 $h$ 的定义）序列 $(\widetilde{J}_{n})$ \trConvergence{}到 $\widetilde{J}$、$(h_{n})$ \trConvergence{}到 $h$，二者都在 $\cC^{0}_{\mathrm{loc}}$ 中，从而也在 $L^{p}_{\mathrm{loc}}$ 中。

现在我们假设已经证明：序列 $(u_{n})$ 在 $W^{k,p}_{\mathrm{loc}}(\RR\times S^{1};\RR^{2n})$ 中\trConvergence{}到 $u$。那么 $(\widetilde{J}_{n})$ 就在 $W^{k,p}_{\mathrm{loc}}(\RR\times S^{1};\RR^{2n})$ 中\trConvergence{}到 $\widetilde{J}$，$(h_{n})$ \trConvergence{}到 $h$。这是如下事实的直接推论：当 $p>2$ 时，$W^{1,p}_{\mathrm{loc}}$ 中的\trConvergence{}蕴含 $\cC^{0}_{\mathrm{loc}}$ 中的\trConvergence。

我们要证明 $(u_{n})$ 在 $W^{k+1,p}_{\mathrm{loc}}(\RR\times S^{1};\RR^{2n})$ 中也\trConvergence{}到 $u$。设 $V\subset U\subset\RR\times S^{1}$ 是相对紧的\trOpen{}子集。从方程
\[
\frac{\partial u}{\partial s}+\widetilde{J}\frac{\partial u}{\partial t}+h=0\qquad\text{与}\qquad\frac{\partial u_{n}}{\partial s}+\widetilde{J}_{n}\frac{\partial u_{n}}{\partial t}+h_{n}=0
\]
我们推出
\[
\frac{\partial(u_{n}-u)}{\partial s}+\widetilde{J}_{n}\frac{\partial(u_{n}-u)}{\partial t}+h_{n}-h+(\widetilde{J}-\widetilde{J}_{n})\frac{\partial u}{\partial t}=0\text{。}
\]
由于 $\widetilde{J}_{n}$ 在 $W^{k,p}_{\mathrm{loc}}(\RR\times S^{1};\RR^{2n})$ 中趋于 $\widetilde{J}$，存在一个 $c_{0}>0$，使得
\[
\|\widetilde{J}_{n}\|_{W^{k,p}(U)}\le c_{0}\qquad\text{对每个 }n\in\NN\text{。}
\]
由\cref{thm:smooth-solution-wkp-estimate}，存在一个不依赖于 $n$ 的常数 $c>0$，使得
\[
\begin{aligned}
\|u_{n}-u\|_{W^{k+1,p}(V)}\le c\bigl(&\|h_{n}-h\|_{W^{k,p}(U)}+\|u_{n}-u\|_{W^{k,p}(U)}\\
&+\bigl\|(\widetilde{J}_{n}-\widetilde{J})\frac{\partial u}{\partial t}\bigr\|_{W^{k,p}(U)}\bigr)\text{。}
\end{aligned}
\]
由于 $\widetilde{J}_{n}$ 在 $W^{k,p}(\overline{U})$ 中趋于 $\widetilde{J}$、$h_{n}$ 趋于 $h$，而 $u$ 是 $\cC^{\infty}$ 类的，我们推出 $u_{n}$ 在 $W^{k+1,p}(V)$ 中趋于 $u$。按归纳法，$u_{n}$ 从而对每个 $k$ 都在 $W^{k,p}_{\mathrm{loc}}$ 中趋于 $u$，因而也在 $\cC^{\infty}_{\mathrm{loc}}$ 中趋于 $u$（仍旧由\trRellich{}定理即\cref{thm:sobolev-estimates}）。

\subsection{如何由定理~12.1.5 与命题~12.1.4 推出引理~12.1.1}
\label{subsec:deduce-lemma-1211}

鉴于\cref{thm:loc-topologies-coincide}，只需证明 $(u_{n})$ 的一个子序列在 $W^{1,p}_{\mathrm{loc}}$ \trTopology{}下有极限。由\trArzelaAscoliTheorem，存在一个子序列——我们仍把它记作 $(u_{n})$——它在 $\cC^{0}_{\mathrm{loc}}$ 中\trConvergence{}到一个极限 $u$。如今只需证明下面这条。

\begin{lemma}\label{thm:un-cauchy-w1p}
对 $U$ 的每个相对紧\trOpen{}子集 $V$，序列 $(u_{n})$ 都是 $W^{1,p}_{\mathrm{loc}}(V)$ \trTopology{}下的\trCauchy{}序列。
\end{lemma}

\begin{proof}
我们想把\cref{thm:generalized-regularity-estimate}的那条不等式\footnote{译注：原书此处印作 “the equality of Theorem 12.1.5”（定理 12.1.5 的等式）；本证明所应用的实为该定理的不等式（下文即有 “The inequality of Theorem 12.1.5 can be written as” 之语），疑应为 “the inequality”。此处径改。——译者注}应用到 $u_{n}-u_{m}$ 上，取 $\ell=0$。由于这是一个局部性的陈述，我们可以、而且也确实要假设每个 $u_{m}$（定义在 $V$ 上并且）在\trEuclidean{}空间 $\RR^{2n}$ 中取值。让我们来验证\cref{thm:generalized-regularity-estimate}的条件得到满足。对 $m,n\in\NN$，我们有
\[
\begin{aligned}
&\frac{\partial(u_{n}-u_{m})}{\partial s}+J(u_{n})\frac{\partial(u_{n}-u_{m})}{\partial t}\\
&\qquad=(J(u_{m})-J(u_{n}))\frac{\partial u_{m}}{\partial t}+\operatorname{grad}H_{t}(u_{m})-\operatorname{grad}H_{t}(u_{n})\text{。}
\end{aligned}
\]
令 $\widetilde{J}=J(u_{n})$，并令
\[
h=(J(u_{m})-J(u_{n}))\frac{\partial u_{m}}{\partial t}+\operatorname{grad}H_{t}(u_{m})-\operatorname{grad}H_{t}(u_{n})\text{。}
\]
我们先对 $k=0$ 验证\cref{thm:generalized-regularity-estimate}的条件。首先，$u_{n}$ 与 $J$ 都是 $\cC^{1}$ 类的，故 $\widetilde{J}\in W^{1,p}(V)$。其次，
\[
\begin{aligned}
\left\|\frac{\partial J(u_{n})}{\partial s}\right\|_{L^{p}(V)}&=\left\|dJ(u_{n})\cdot\frac{\partial u_{n}}{\partial s}\right\|_{L^{p}(V)}\\
&\le\sup_{w\in W}\|(dJ)_{w}\|\cdot\left\|\frac{\partial u_{n}}{\partial s}\right\|_{L^{p}(V)}\\
&\le\sup_{w\in W}\|(dJ)_{w}\|\cdot M\cdot(\operatorname{area}(V))^{1/p}\text{，}
\end{aligned}
\]
这个上界不依赖于 $n$。同样的上界对 $\|\partial J(u_{n})/\partial t\|_{L^{p}(V)}$ 与
\[
\|J(u_{n})\|_{L^{p}(V)}\le\sup_{w\in W}\|J_{w}\|\cdot(\operatorname{area}(V))^{1/p}
\]
也成立。我们于是发现 $\|\widetilde{J}\|_{W^{1,p}(V)}\le c_{0}$。此时 $h$ 是连续的，因而是 $L^{p}(V)$ 的一个元素。于是我们可以对 $k=1$、$\ell=0$ 应用\cref{thm:generalized-regularity-estimate}：设 $V'\supset V$ 是 $U$ 的一个相对紧\trOpen{}子集。\cref{thm:generalized-regularity-estimate}的那条不等式可以写成
\[
\|u_{n}-u_{m}\|_{W^{1,p}(V)}\le C\left(\|u_{n}-u_{m}\|_{L^{p}(V')}+\|h\|_{L^{p}(V')}\right)\text{。}
\]
我们可以把右边的第二项从上方加以估计：
\[
\begin{aligned}
\|h\|_{L^{p}(V')}&\le\left\|(J(u_{n})-J(u_{m}))\frac{\partial u_{m}}{\partial t}\right\|_{L^{p}(V')}\\
&\qquad+\|\operatorname{grad}H_{t}(u_{m})-\operatorname{grad}H_{t}(u_{n})\|_{L^{p}(V')}\\
&\le\sup_{w\in W}\|(dJ)_{w}\|\cdot\|u_{n}-u_{m}\|_{L^{p}(V')}\cdot M\\
&\qquad+\sup_{W\times S^{1}}\|d(\operatorname{grad}H)\|\cdot\|u_{n}-u_{m}\|_{L^{p}(V')}\\
&\le K\cdot\|u_{n}-u_{m}\|_{L^{p}(V')}\text{。}
\end{aligned}
\]
此时，该序列在 $\cC^{0}_{\mathrm{loc}}$ 中是\trConvergence{}的，故在 $L^{p}(V')$ 中\trConvergence，从而在后者中是\trCauchy{}序列。这蕴含它在 $W^{1,p}$ 中也是\trCauchy{}序列，引理得证。
\end{proof}

至此，借助于\cref{thm:generalized-regularity-estimate}（经由\cref{thm:loc-topologies-coincide}），我们已经证明了\trRegularity{}引理，即\cref{thm:elliptic-regularity-lemma}。这条定理本身仍有待证明，我们下面就来做这件事。

\subsection{定理~12.1.5 的证明}
\label{subsec:proof-thm-1215}

当 $\widetilde{J}$ 就是（常值且等于）标准\trComplexStructure{} $J_{0}$ 时，我们认出的正是\cref{thm:local-regularity-estimate}的陈述。在一般的情形下，证明更为复杂，而且条件 $p>2$ 是本质性的（注意，要证明\cref{thm:local-regularity-estimate}，$p>1$ 便已足够）。

\emph{第一步，$k=0$ 的情形。}

我们先对 $k=0$（且 $\ell=0$）证明这条定理。沿着\cref{thm:local-regularity-estimate}的证明的策略，我们很快便遇到一个困难：把算子 $\partial/\partial s-\widetilde{J}\,\partial/\partial t$ 应用到方程
\[
\frac{\partial u}{\partial s}+\widetilde{J}\frac{\partial u}{\partial t}+h=0
\]
上时，我们得到\footnote{译注：下式第二、三项原书均印作 $\cdot\,u$（乘的是 $u$）；按乘积法则，把 $\partial/\partial s-\widetilde{J}\,\partial/\partial t$ 作用于 $\widetilde{J}\,\partial u/\partial t$ 所出的这两项应是 $(\partial\widetilde{J}/\partial s)\cdot(\partial u/\partial t)$ 与 $-\widetilde{J}(\partial\widetilde{J}/\partial t)\cdot(\partial u/\partial t)$——下文引理~12.4.3 证明中对 $J$ 的同款展开（该处即印作 $\partial J/\partial s\cdot\partial u/\partial t$ 之形，无误）可资对照，属局部公式之误。照排。——译者注}
\[
\Delta u+\frac{\partial\widetilde{J}}{\partial s}\cdot u-\widetilde{J}\frac{\partial\widetilde{J}}{\partial t}\cdot u+\frac{\partial h}{\partial s}-\widetilde{J}\frac{\partial h}{\partial t}=0\text{。}
\]
我们无法应用\cref{thm:bar-partial-regularity-criterion}，因为作为诸 $L^{p}_{\mathrm{loc}}$ 函数的一个乘积，
\[
\frac{\partial}{\partial s}\widetilde{J}\cdot u\in L^{p/2}_{\mathrm{loc}}
\]
（仅仅！）由\trHoelder{}不等式可得。为了能应用同一个方法，我们先把\cref{thm:generalized-regularity-estimate}的陈述修改如下。

\begin{lemma}\label{thm:nonlinear-regularity-triple}
设 $U$ 与 $V$ 是 $\RR\times S^{1}$ 的两个满足 $\overline{V}\subset U$ 的相对紧\trOpen{}子集，并设 $p$、$q$、$r\in\RR^{+}$ 满足
\[
p>2\text{，}\qquad r>1\qquad\text{且}\qquad\frac{1}{p}+\frac{1}{q}=\frac{1}{r}\text{。}
\]
设 $u$ 是方程
\[
\frac{\partial u}{\partial s}+J\frac{\partial u}{\partial t}=f
\]
（在\trDistribution{}的意义下的）一个解（其中 $J$ 是 $W^{1,p}\bigl(U;\operatorname{End}(\RR^{2n})\bigr)$ 中的一个\trAlmostComplex{}结构，满足 $\|J\|_{W^{1,p}(U)}\le c_{0}$，而 $f\in L^{r}_{\mathrm{loc}}(U;\RR^{2n})$）。那么：
（1）若 $u\in L^{q}_{\mathrm{loc}}(U;\RR^{2n})$，则 $u\in W^{1,r}_{\mathrm{loc}}(U;\RR^{2n})$。

（2）存在一个依赖于 $c_{0}$、$U$ 与 $V$ 的常数 $c>0$，使得
\[
\|u\|_{W^{1,r}(V)}\le c\left(\|f\|_{L^{r}(U)}+\|u\|_{L^{q}(U)}\right)\text{。}
\]
\end{lemma}

\begin{remark}\label{thm:nonlinear-regularity-remarks}
（1）若 $r<2$，那么 $p$、$q$ 与 $r$ 之间的关系便蕴含 $q<2r/(2-r)$。在这种情形下，\cref{thm:rellich-kondrachov}断言 $W^{1,r}_{\mathrm{loc}}\subset L^{q}_{\mathrm{loc}}$。

（2）在\cref{thm:nonlinear-regularity-triple}的陈述中，$p$ 与 $q$ 可以取 $+\infty$ 为值：
\begin{itemize}
\item 若 $p=+\infty$，则 $r=q<+\infty$。
\item 若 $q=+\infty$，则 $r=p<+\infty$。
\end{itemize}
\end{remark}

\begin{proof}[\cref{thm:nonlinear-regularity-triple}的证明]
对引理的方程施加 $\partial/\partial s-J\,\partial/\partial t$，我们便得到
\[
\begin{aligned}
\Delta u&=-\frac{\partial J}{\partial s}\frac{\partial u}{\partial t}+J\frac{\partial J}{\partial t}\frac{\partial u}{\partial t}+\frac{\partial f}{\partial s}-J\frac{\partial f}{\partial t}\\
&=-\frac{\partial J}{\partial s}\frac{\partial u}{\partial t}-\frac{\partial J}{\partial t}J\frac{\partial u}{\partial t}+\frac{\partial f}{\partial s}-J\frac{\partial f}{\partial t}\\
&=-\frac{\partial J}{\partial s}\frac{\partial u}{\partial t}-\frac{\partial J}{\partial t}\left(f-\frac{\partial u}{\partial s}\right)+\frac{\partial f}{\partial s}-J\frac{\partial f}{\partial t}
\end{aligned}
\]
（这里用到了 $J^{2}=-\mathrm{Id}$ 与 $u$ 是一个解这两个事实）。这便引出
\[
\begin{aligned}
\Delta u&=-\frac{\partial J}{\partial s}\frac{\partial u}{\partial t}+\frac{\partial J}{\partial t}\frac{\partial u}{\partial s}+\frac{\partial f}{\partial s}-\frac{\partial}{\partial t}(Jf)\\
&=-\frac{\partial}{\partial t}\left(\frac{\partial J}{\partial s}\cdot u\right)+\frac{\partial^{2}J}{\partial t\partial s}\cdot u+\frac{\partial}{\partial s}\left(\frac{\partial J}{\partial t}u\right)-\frac{\partial^{2}J}{\partial s\partial t}\cdot u+\frac{\partial f}{\partial s}-\frac{\partial}{\partial t}(Jf)\\
&=\frac{\partial}{\partial s}\left(\frac{\partial J}{\partial t}\cdot u+f\right)+\frac{\partial}{\partial t}\left(-\frac{\partial J}{\partial s}\cdot u-Jf\right)\text{。}
\end{aligned}
\]
利用\trHoelder{}不等式与 $J$ 的连续性，我们得到
\[
\frac{\partial J}{\partial t}\cdot u+f\qquad\text{与}\qquad-\frac{\partial J}{\partial s}\cdot u-Jf\in L^{r}_{\mathrm{loc}}(U)
\]
（这对 $p$ 或 $q$ 为无穷的情形也适用）。于是\cref{thm:bar-partial-regularity-criterion}的条件得到满足，这使我们得以应用它的结论，得到一个常数 $c_{1}>0$，使得
\[
\begin{aligned}
\|u\|_{W^{1,r}(V)}&\le c_{1}\left(\left\|\frac{\partial J}{\partial t}\cdot u\right\|_{L^{r}(U)}+\|f\|_{L^{r}(U)}+\left\|\frac{\partial J}{\partial s}\cdot u\right\|_{L^{r}(U)}+\|Jf\|_{L^{r}(U)}\right)\\
&\le\Bigl(\|J\|_{W^{1,p}(U)}\|u\|_{L^{q}(U)}+\|f\|_{L^{r}(U)}\\
&\qquad+\|J\|_{W^{1,p}(U)}\|u\|_{L^{q}(U)}+\|J\|_{L^{\infty}(U)}\|f\|_{L^{r}(U)}\Bigr)\\
&\le c_{1}\bigl(2c_{0}\|u\|_{L^{q}(U)}+(1+kc_{0})\|f\|_{L^{r}(U)}\bigr)
\end{aligned}
\]
（其中我们用到了包含关系 $W^{1,p}(U)\subset L^{\infty}(U)$ 的连续性）。\footnote{译注：上式中 $(1+kc_{0})$ 的 $k$ 在本证明中没有定义；按上一行 $\|Jf\|_{L^{r}(U)}\le\|J\|_{L^{\infty}(U)}\|f\|_{L^{r}(U)}$ 与这一句所援引的连续\trEmbedding——$\|J\|_{L^{\infty}(U)}\le K\|J\|_{W^{1,p}(U)}\le Kc_{0}$——$k$ 应即\trSobolev\trEmbedding{}中的常数 $K$（当然亦可径把它吸入下文的新常数 $c$）。照排。——译者注}由此我们推出，存在一个（不依赖于 $u$ 与 $J$ 的）常数 $c$，使得
\[
\|u\|_{W^{1,r}(V)}\le c\left(\|f\|_{L^{r}(U)}+\|u\|_{L^{q}(U)}\right)\text{，}
\]
引理得证。
\end{proof}

那么，让我们回到\cref{thm:generalized-regularity-estimate}的 $k=0$ 情形的证明。

我们先证明 $u\in W^{1,p}_{\mathrm{loc}}(\RR\times S^{1})$。直接应用\cref{thm:nonlinear-regularity-triple}只能给出 $u\in W^{1,p/2}_{\mathrm{loc}}(\RR\times S^{1})$。\footnote{译注：原书本句印作 $u\in W^{1,p/2}(\RR\times S^{1})$（无 $\mathrm{loc}$）；上文同页同款 $W^{2,p}$ 一处原书亦然（彼处已注），此处一并径改补 $\mathrm{loc}$。——译者注}为了得到所期望的\trRegularity，我们需要把这个引理应用到 $q=\infty$ 的情形，也就是说，要带着条件 $u\in L^{\infty}_{\mathrm{loc}}(\RR\times S^{1})$——而只要对某个 $r>2$ 有 $u\in W^{1,r}_{\mathrm{loc}}(\RR\times S^{1})$，这条条件便成立。数 $r=p/2$ 并不满足这条不等式，但一个自举式的论证使我们得以同时提高 $u$ 的\trRegularity{}与 $r$ 的值。

我们假设已经证明：对某个 $r\in\,]1,2[$ 有 $u\in W^{1,r}_{\mathrm{loc}}(\RR\times S^{1})$。利用\cref{thm:rellich-kondrachov}的\trSobolev\trEmbedding，我们可以推出：对 $q=2r/(2-r)$ 有 $u\in L^{q}_{\mathrm{loc}}(\RR\times S^{1})$。随后，由\cref{thm:nonlinear-regularity-triple}，我们有 $u\in W^{1,r'}_{\mathrm{loc}}(\RR\times S^{1})$，其中
\[
\frac{1}{r'}=\frac{1}{p}+\frac{1}{q}\text{，}\qquad\text{也就是说}\qquad r'=\frac{pq}{p+q}=\frac{2pr}{2p+2r-pr}
\]
（$h$ 属于 $L^{p}_{\mathrm{loc}}$，从而特别地，当 $1<r'\le p$ 时属于 $L^{r'}_{\mathrm{loc}}$）。我们看到，若 $p>2$，则 $r'>r$。\emph{条件 $p>2$ 因而对提高 $u$ 的\trRegularity{}是本质性的。}于是我们考虑满足下述条件的序列 $(r_{m})$：
\[
r_{0}\in\,]1,2[\text{，}\qquad r_{0}\le\frac{p}{2}\qquad\text{与}\qquad r_{m+1}=\frac{2pr_{m}}{2p+2r_{m}-pr_{m}}\text{，}
\]
于是，若 $r_{m}<2$，则 $r_{m+1}>r_{m}$，并且 $u\in W^{1,r_{m+1}}_{\mathrm{loc}}(\RR\times S^{1})$。若对每个 $m$ 都有 $r_{m}\in\,]1,2[$，那么递增序列 $r_{m}$ 便\trConvergence{}到一个极限 $\ell\in\,]1,2]$……而由于 $p>2$，这是递推关系所不允许的。因而我们对某个 $m$ 有 $r_{m}<2\le r_{m+1}$，从而对每个满足 $1<r\le 2$ 的 $r$ 都有 $u\in W^{1,r}_{\mathrm{loc}}(\RR\times S^{1})$。另一方面，
\[
r'-2=\frac{2pr}{2p+2r-pr}-2=\frac{4[(p-1)(r-1)-1]}{2p+2r-pr}
\]
当 $r\in\,]1,2]$ 充分接近于 $2$ 时是正的。由此可见
\[
u\in W^{1,r'}_{\mathrm{loc}}(\RR\times S^{1})\subset L^{\infty}_{\mathrm{loc}}(\RR\times S^{1})\qquad\text{从而}\qquad u\in W^{1,p}_{\mathrm{loc}}(\RR\times S^{1})\text{，}
\]
这最后一步由\cref{thm:nonlinear-regularity-triple}得到。

我们现在来证明不等式
\[
\|u\|_{W^{1,p}(V)}\le c\left(\|h\|_{L^{p}(U)}+\|u\|_{L^{p}(U)}\right)\text{。}
\]
为了证明 $u\in W^{1,p}_{\mathrm{loc}}(\RR\times S^{1})$，我们对诸三元组
\[
(p,q_{0},r_{0})\text{，}\ (p,q_{1},r_{1})\text{，}\ \dots\text{，}\ (p,q_{m},r_{m})\text{，}\ (p,q_{m+1},r_{m+1})\text{，}\ (p,+\infty,p)
\]
相继应用\cref{thm:nonlinear-regularity-triple}，它们满足下列性质：
\begin{itemize}
\item 诸 $r_{m}$ 递增，更确切地说，
\[
1<r_{0}<r_{1}<\dots<r_{m}<2<r_{m+1}<p\text{。}
\]
\item 此外，$q_{0}=p$，并且对每个 $i\ge 1$ 有 $q_{i}=2r_{i-1}/(2-r_{i-1})$。\footnote{译注：这组性质与\cref{thm:nonlinear-regularity-triple}的前提不\trCompatible：该引理要求 $1/p+1/q_{0}=1/r_{0}$，取 $q_{0}=p$ 便强制 $r_{0}=p/2$，而上文（原书 p.~470）的序列只设 $r_{0}\in\,]1,2[$、$r_{0}\le p/2$——$r_{0}<p/2$ 时违反该关系式，$p\ge 4$ 时 $r_{0}=p/2$ 又不在 $]1,2[$ 内；下文“第一次应用”（对 $(p,q_{0},r_{0})=(p,p,r_{0})$）以及第二步开头（原书 p.~473，$q_{0}=p$）的应用皆受同一问题影响。修补可照此进行：取 $r_{0}\in\,]1,\min(2,p/2)]$，按关系式定 $q_{0}$，即 $1/q_{0}=1/r_{0}-1/p$（从而 $q_{0}=pr_{0}/(p-r_{0})\le p$）；此时局部地（$\overline{V}$ 为\trCompact{}时）$L^{p}\subset L^{q_{0}}$，故由 $u\in L^{p}_{\mathrm{loc}}$ 得 $u\in L^{q_{0}}_{\mathrm{loc}}$，首次应用的结论遂为 $\|u\|_{W^{1,r_{0}}(V_{0})}\le c\bigl(\|h\|_{L^{p}(U)}+\|u\|_{L^{q_{0}}(U)}\bigr)$，再归入 $\|u\|_{W^{1,r_{0}}(V_{0})}\le c\bigl(\|h\|_{L^{p}(U)}+\|u\|_{L^{p}(U)}\bigr)$，终式不变。亦可按 $p$ 的区间分别处理：$2<p<4$ 时取 $r_{0}=p/2$、$q_{0}=p$ 即自洽。照排。——译者注}
\end{itemize}
\cref{thm:nonlinear-regularity-triple}的每一次应用都给出 $W^{1,r}$、$L^{r}$ 与 $L^{q}$ 诸范数之间的一条不等式。为了写出这些不等式，我们选取相对紧的\trOpen{}集 $(V_{i})_{i=0,\dots,m}$，使得
\[
V\subset V_{m+1}\subset V_{m}\subset\dots\subset V_{0}\subset U\text{。}
\]
\cref{thm:nonlinear-regularity-triple}的最后一次应用，即对三元组 $(p,+\infty,p)$ 的应用，给出
\[
\begin{aligned}
\|u\|_{W^{1,p}(V)}&\le k_{1}\left(\|h\|_{L^{p}(V_{m+1})}+\|u\|_{L^{\infty}(V_{m+1})}\right)\\
&\le k_{2}\left(\|h\|_{L^{p}(U)}+\|u\|_{W^{1,r_{m+1}}(V_{m+1})}\right)
\end{aligned}
\]
带有连续\trEmbedding
\[
W^{1,r_{m+1}}(V_{m+1})\hookrightarrow L^{\infty}(V_{m+1})\text{，}\qquad\text{因为 }r_{m+1}>2\text{。}
\]
其次，对 $i\in\{1,\dots,m+1\}$，我们有
\[
\begin{aligned}
\|u\|_{W^{1,r_{i}}(V_{i})}&\le c_{i}\left(\|h\|_{L^{r_{i}}(V_{i-1})}+\|u\|_{L^{q_{i}}(V_{i-1})}\right)\\
&\le C_{i}\left(\|h\|_{L^{p}(V_{i-1})}+\|u\|_{W^{1,r_{i-1}}(V_{i-1})}\right)
\end{aligned}
\]
带有连续\trEmbedding
\[
L^{p}(V_{i-1})\subset L^{r_{i}}(V_{i-1})\qquad\text{当 }r_{i}<p\text{ 且 }\overline{V}_{i-1}\text{ 为\trCompact{}时}
\]
以及
\[
W^{1,r_{i-1}}(V_{i-1})\subset L^{q_{i}}(V_{i-1})\qquad\text{当 }q_{i}=\frac{2r_{i-1}}{2-r_{i-1}}\text{ 且 }\overline{V}_{i-1}\text{ 为\trCompact{}时。}
\]
从而，对每个 $i=1,\dots,m+1$，我们有
\[
\|u\|_{W^{1,r_{i}}(V_{i})}\le C_{i}\left(\|h\|_{L^{p}(U)}+\|u\|_{W^{1,r_{i-1}}(V_{i-1})}\right)\text{。}
\]
把所有这些不等式结合起来，我们发现
\[
\|u\|_{W^{1,r_{m+1}}(V_{m+1})}\le k_{3}\left(\|h\|_{L^{p}(U)}+\|u\|_{W^{1,r_{0}}(V_{0})}\right)\text{，}
\]
它与先前得到的不等式
\[
\|u\|_{W^{1,p}(V)}\le k_{2}\left(\|h\|_{L^{p}(U)}+\|u\|_{W^{1,r_{m+1}}(V_{m+1})}\right)
\]
一起便给出
\[
\|u\|_{W^{1,p}(V)}\le k_{4}\left(\|h\|_{L^{p}(U)}+\|u\|_{W^{1,r_{0}}(V_{0})}\right)\text{。}
\]
最后，\cref{thm:nonlinear-regularity-triple}的第一次应用，即对 $(p,q_{0},r_{0})=(p,p,r_{0})$ 的应用，给出
\[
\|u\|_{W^{1,r_{0}}(V_{0})}\le c_{0}\left(\|h\|_{L^{p}(U)}+\|u\|_{L^{p}(U)}\right)\text{，}
\]
结合前面的不等式，这就蕴含\footnote{译注：原书此处左端印作 $\|u\|_{W^{1,p}(U)}$；按本页上文所列的证明目标不等式（左端为 $\|u\|_{W^{1,p}(V)}$）及整条估计链（诸上界皆对 $V$ 而言），疑应为 $\|u\|_{W^{1,p}(V)}$。此处径改。——译者注}
\[
\|u\|_{W^{1,p}(V)}\le c\left(\|h\|_{L^{p}(U)}+\|u\|_{L^{p}(U)}\right)\text{，}
\]
即我们所要证的那条不等式。

\emph{第二步，$k=1$ 的情形。}

让我们来证明 $u\in W^{2,p}_{\mathrm{loc}}(\RR\times S^{1})$。为简单起见，我们将写 $v=\partial u/\partial s$（以及 $w=\partial u/\partial t$）。那么 $v$ 是方程
\[
\frac{\partial v}{\partial s}+\widetilde{J}\frac{\partial v}{\partial t}+g=0
\]
的一个解，其中我们已令
\[
g(s,t)=\frac{\partial h}{\partial s}(s,t)+\frac{\partial\widetilde{J}}{\partial s}(s,t)\frac{\partial u}{\partial t}(s,t)\text{。}
\]
这里的情形与第一步中所考虑的有所不同，因为此处出现的 $g$ 只属于 $L^{p/2}_{\mathrm{loc}}(\RR\times S^{1})$（仍旧由\trHoelder{}不等式），而不再属于 $L^{p}_{\mathrm{loc}}(\RR\times S^{1})$。我们可以把\cref{thm:nonlinear-regularity-triple}应用到三元组 $(p,q_{0},r_{0})$ 上——取 $q_{0}=p$，且像前一个证明中那样取 $r_{0}\in\,]1,2[$、$r_{0}\le p/2$——\footnote{译注：原书此处印作 $r_{0}=p/2$；按第一步中的取法 $r_{0}\in\,]1,2[$、$r_{0}\le p/2$（$=p/2$ 仅当 $p<4$ 时才落入 $]1,2[$，而原书只设 $p>2$），疑应为 $r_{0}\le p/2$。此处径改。——译者注}并得出 $v\in W^{1,r_{0}}_{\mathrm{loc}}(\RR\times S^{1})$。以类似的方式，我们可以证明 $w\in W^{1,r_{0}}_{\mathrm{loc}}(\RR\times S^{1})$。

像上面那样继续下去，我们得到：对 $q_{1}=2r_{0}/(2-r_{0})$，有 $v,w\in L^{q_{1}}_{\mathrm{loc}}(\RR\times S^{1})$（这正是\trSobolev\trEmbedding）。于是，由\trHoelder{}不等式，
\[
\frac{\partial\widetilde{J}}{\partial s}\cdot\frac{\partial u}{\partial t}\quad\text{与}\quad\frac{\partial\widetilde{J}}{\partial t}\cdot\frac{\partial u}{\partial t}\in L^{r_{1}}_{\mathrm{loc}}(\RR\times S^{1})\text{，}\qquad\text{其中 }\frac{1}{p}+\frac{1}{q_{1}}=\frac{1}{r_{1}}\text{。}
\]
从而 $g\in L^{r_{1}}_{\mathrm{loc}}$。于是我们可以对三元组 $(p,q_{1},r_{1})$ 再一次应用\cref{thm:nonlinear-regularity-triple}，并发现 $v$ 与 $w$ 都属于 $W^{1,r_{1}}_{\mathrm{loc}}(\RR\times S^{1})$。

于是我们可以像先前那样相继应用\cref{thm:nonlinear-regularity-triple}，因为每一次应用都改进了 $g$ 的可积性（即 $g$ 所属于的那个 $L^{r}$ 中的 $r$）。最后，我们得到
\[
v=\frac{\partial u}{\partial s}\qquad\text{与}\qquad w=\frac{\partial u}{\partial t}\in W^{1,p}_{\mathrm{loc}}(\RR\times S^{1})\text{，}
\]
换言之，$u\in W^{2,p}_{\mathrm{loc}}(\RR\times S^{1})$，\footnote{译注：原书此处印作 $u\in W^{2,p}(\mathbf{R}\times S^{1})$（无 loc）；按本步开头“Let us show that $u\in W^{2,p}_{\mathrm{loc}}(\mathbf{R}\times S^{1})$”的目标句及全证的局部性，疑为漏排，应为 $W^{2,p}_{\mathrm{loc}}(\mathbf{R}\times S^{1})$。此处径改（补 $\mathrm{loc}$）。——译者注}这正是我们所要证明的。这就结束了第二步，即 $k=1$ 的情形。

\par\nopagebreak
\emph{第三步，对 $k$ 的归纳。}\nopagebreak

现在让我们假设\cref{thm:generalized-regularity-estimate}已经对 $k$ 得证，并来对 $k+1$ 证明它。设\footnote{译注：下式末款原书印作 $\widetilde{h}\in W^{1,k+1}_{\mathrm{loc}}(\mathbf{R}\times S^{1};\mathbf{R}^{2n})$。其一，本归纳所设的\cref{thm:generalized-regularity-estimate}中该数据恒记作 $h$（紧随其后的方程亦用 $h$），波纹号疑为衍字，应为 $h$；其二，按归纳假设该数据应属于 $W^{k+1,p}_{\mathrm{loc}}$（与同句 $\widetilde{J}\in W^{k+1,p}_{\mathrm{loc}}$ 平行），上标 $1,k+1$ 疑为 $k+1,p$ 之误。此处均径改。——译者注}
\[
\widetilde{J}\in W^{k+1,p}_{\mathrm{loc}}\bigl(\RR\times S^{1};\operatorname{End}(\RR^{2n})\bigr)\text{，}\ \widetilde{J}^{2}=-\mathrm{Id}\text{，}\ h\in W^{k+1,p}_{\mathrm{loc}}(\RR\times S^{1};\RR^{2n})\text{，}
\]
设 $u\in L^{p}_{\mathrm{loc}}(\RR\times S^{1};\RR^{2n})$ 是
\[
\frac{\partial u}{\partial s}+\widetilde{J}\frac{\partial u}{\partial t}+h=0
\]
的一个\trWeakSolution。按归纳假设，$u\in W^{k+1,p}_{\mathrm{loc}}(\RR\times S^{1};\RR^{2n})$。我们要证明 $u\in W^{k+2,p}_{\mathrm{loc}}(\RR\times S^{1};\RR^{2n})$。像上面那样，我们把它的导数记为
\[
v=\frac{\partial u}{\partial s}\qquad\text{与}\qquad w=\frac{\partial u}{\partial t}\text{。}
\]
我们已经看到，$v$ 满足
\[
\frac{\partial v}{\partial s}+\widetilde{J}\frac{\partial v}{\partial t}+\frac{\partial\widetilde{J}}{\partial s}\cdot\frac{\partial u}{\partial t}+\frac{\partial h}{\partial s}=0\text{。}
\]
为了能把归纳假设应用于它，我们需要知道
\[
\frac{\partial\widetilde{J}}{\partial s}\cdot\frac{\partial u}{\partial t}\in W^{k,p}_{\mathrm{loc}}(\RR\times S^{1})\text{。}
\]
我们将由下面这条更一般的结果得到这一点。

\begin{lemma}\label{thm:product-wkp-estimate}
设 $p>2$，并设 $k\ge 1$。若 $f$ 与 $g$ 都属于 $W^{k,p}_{\mathrm{loc}}(\RR\times S^{1};\RR)$，那么 $fg\in W^{k,p}_{\mathrm{loc}}(\RR\times S^{1};\RR)$，而且对每个相对紧的\trOpen{}子集 $V$，都存在一个常数 $K>0$，使得
\[
\|fg\|_{W^{k,p}(V)}\le K\|f\|_{W^{k,p}(V)}\|g\|_{W^{k,p}(V)}\text{。}
\]
\end{lemma}

\begin{proof}
这基本上是莱布尼茨法则的一个推论。设 $V$ 是一个相对紧的\trOpen{}集。我们按归纳法进行。首先令 $k=1$。我们利用连续\trEmbedding{} $W^{1,p}(V)\subset L^{\infty}(V)$。\footnote{译注：$V$ 仅是相对紧的\trOpen{}集，而 $W^{1,p}(V)\subset L^{\infty}(V)$ 这类\trSobolev\trEmbedding{}的现成陈述通常以区域边界\trRegular{}（如 \trLipschitz{}域）或可延拓为前提，对任意相对紧开集直接援引需要补充论证。本章的用法可照此补救：在 $\overline{V}$ 之外再取一个（或有限个）\trSmooth{}小域作内估计，\trEmbedding{}常数随集而变，不影响论证；下文 $\ell=1$ 处对 $U$ 的同类用法 $\bigl\|\partial\widetilde{J}/\partial s\bigr\|_{L^{\infty}(U)}\le K\bigl\|\partial\widetilde{J}/\partial s\bigr\|_{W^{1,p}(U)}$ 亦同此补（或径设所论集合为 \trLipschitz{}域／可延拓域）。照排。——译者注}我们有
\[
\begin{aligned}
\|fg\|_{L^{p}(V)}&\le\|f\|_{L^{p}(V)}\|g\|_{L^{\infty}(V)}\\
&\le K_{1}\|f\|_{L^{p}(V)}\|g\|_{W^{1,p}(V)}\\
&\le K_{1}\|f\|_{W^{1,p}(V)}\|g\|_{W^{1,p}(V)}\text{，}
\end{aligned}
\]
\[
\begin{aligned}
\left\|\frac{\partial(fg)}{\partial s}\right\|_{L^{p}(V)}&\le\left\|\frac{\partial f}{\partial s}g\right\|_{L^{p}(V)}+\left\|f\frac{\partial g}{\partial s}\right\|_{L^{p}(V)}\\
&\le\left\|\frac{\partial f}{\partial s}\right\|_{L^{p}(V)}\|g\|_{L^{\infty}(V)}+\left\|\frac{\partial g}{\partial s}\right\|_{L^{p}(V)}\|f\|_{L^{\infty}(V)}\\
&\le K_{2}\|f\|_{W^{1,p}(V)}\|g\|_{W^{1,p}(V)}\text{，}
\end{aligned}
\]
对 $fg$ 关于 $t$ 的导数也有一个类似的不等式。于是 $k=1$ 的情形得证。我们假设所述命题对 $k$ 成立。我们有
\[
\begin{aligned}
\left\|\frac{\partial(fg)}{\partial s}\right\|_{W^{k,p}(V)}&\le\left\|\frac{\partial f}{\partial s}g\right\|_{W^{k,p}(V)}+\left\|f\frac{\partial g}{\partial s}\right\|_{W^{k,p}(V)}\\
&\le K_{3}\left\|\frac{\partial f}{\partial s}\right\|_{W^{k,p}(V)}\|g\|_{W^{k,p}(V)}+K_{4}\|f\|_{W^{k,p}(V)}\left\|\frac{\partial g}{\partial s}\right\|_{W^{k,p}(V)}\\
&\le K\|f\|_{W^{k+1,p}(V)}\|g\|_{W^{k+1,p}(V)}
\end{aligned}
\]
这里用到了归纳假设。对 $fg$ 关于 $t$ 的导数也有一个类似的不等式。于是
\[
\|fg\|_{W^{k+1,p}(V)}\le K\|f\|_{W^{k+1,p}(V)}\|g\|_{W^{k+1,p}(V)}
\]
引理得证。
\end{proof}

借助这条引理（即\cref{thm:product-wkp-estimate}），我们如今可以应用归纳假设，并得到
\[
v=\frac{\partial u}{\partial s}\qquad\text{与}\qquad w=\frac{\partial u}{\partial t}\in W^{k+1,p}_{\mathrm{loc}}(\RR\times S^{1};\RR^{2n})\text{，}
\]
从而 $u\in W^{k+2,p}_{\mathrm{loc}}(\RR\times S^{1};\RR^{2n})$。

尚须证明关于范数的那条不等式。设 $U$ 与 $V$ 是满足 $\overline{V}\subset U$ 的相对紧\trOpen{}集，并设 $\ell\in\{0,\dots,k\}$。\footnote{译注：原书此处作 $\ell\in\{0,\dots,k\}$，本身不误；但 §12.1 中\cref{thm:generalized-regularity-estimate}的陈述处原书印作 $\ell\in\{0,\dots,\max(k-1,0)\}$，与本证明的实际用法（此处及下文对每个 $\ell\le k$ 的估计）相抵，彼处的 $\max(k-1,0)$ 系原书之误，应为 $k$。彼处已径改（2026-09-10 外部复审 RM12-001，见该处译注）。——译者注}对 $\ell=0$，所要的不等式已在上面得到证明。对 $\ell\ge 1$，我们把归纳假设应用于 $v$（对 $w$ 亦然），取 $\ell-1$，这就给出
\[
\|v\|_{W^{\ell,p}(V)}\le c\left(\left\|\frac{\partial h}{\partial s}\right\|_{W^{\ell-1,p}(U)}+\left\|\frac{\partial\widetilde{J}}{\partial s}\frac{\partial u}{\partial t}\right\|_{W^{\ell-1,p}(U)}+\|v\|_{W^{\ell-1,p}(U)}\right)\text{。}
\]
对 $\ell\ge 2$，由\cref{thm:product-wkp-estimate}，这条不等式给出
\[
\begin{aligned}
\|v\|_{W^{\ell,p}(V)}&\le c\Bigl(\|h\|_{W^{\ell,p}(U)}+K\left\|\frac{\partial\widetilde{J}}{\partial s}\right\|_{W^{\ell-1,p}(U)}\left\|\frac{\partial u}{\partial t}\right\|_{W^{\ell-1,p}(U)}\\
&\qquad\qquad\qquad\qquad+\|v\|_{W^{\ell-1,p}(U)}\Bigr)\\
&\le\bigl(\|h\|_{W^{\ell,p}(U)}+Kc_{0}\|u\|_{W^{\ell,p}(U)}+\|u\|_{W^{\ell,p}(U)}\bigr)\text{。}
\end{aligned}
\]
对 $\ell=1$，我们利用 $\widetilde{J}\in W^{k+1,p}(V)$（其中 $k\ge\ell=1$）来推出 $\partial\widetilde{J}/\partial s\in W^{1,p}$，于是
\[
\begin{aligned}
\left\|\frac{\partial\widetilde{J}}{\partial s}\frac{\partial u}{\partial t}\right\|_{L^{p}(U)}&\le\left\|\frac{\partial\widetilde{J}}{\partial s}\right\|_{L^{\infty}(U)}\left\|\frac{\partial u}{\partial t}\right\|_{L^{p}(U)}\\
&\le K\left\|\frac{\partial\widetilde{J}}{\partial s}\right\|_{W^{1,p}(U)}\|u\|_{W^{1,p}(U)}\\
&\le K\|\widetilde{J}\|_{W^{2,p}(U)}\|u\|_{W^{1,p}(U)}\\
&\le Kc_{0}\|u\|_{W^{1,p}(U)}
\end{aligned}
\]
这是因为 $\|\widetilde{J}\|_{W^{k+1,p}(U)}\le c_{0}$。从而我们还有
\[
\|v\|_{W^{1,p}(V)}\le c\left(\|h\|_{W^{1,p}(U)}+Kc_{0}\|u\|_{W^{1,p}(U)}+\|u\|_{W^{1,p}(U)}\right)\text{。}
\]
这样，对每个 $\ell=0,\dots,k$（并且对 $w$ 与 $v$ 二者），我们都证明了存在如下形式的一个上界：
\[
\|v\|_{W^{\ell,p}(V)}\le c\left(\|h\|_{W^{\ell,p}(U)}+\|u\|_{W^{\ell,p}(U)}\right)\text{。}
\]
写出
\[
\|u\|_{W^{\ell+1,p}(V)}\le\|u\|_{W^{\ell,p}(V)}+\|v\|_{W^{\ell,p}(V)}+\|w\|_{W^{\ell,p}(V)}
\]
并利用所有这些上界，我们便能证明所要的那条不等式，\cref{thm:generalized-regularity-estimate}的证明也就此结束。
